%%%%%%%%%%%%%%%%%%%%%%%%%%%%%%%%%%%%%%%%%
% Medium Length Professional CV
% LaTeX Template
% Version 3.0 (December 17, 2022)
%
% This template originates from:
% https://www.LaTeXTemplates.com
%
% Author:
% Vel (vel@latextemplates.com)
%
% Original author:
% Trey Hunner (http://www.treyhunner.com/)
%
% License:
% CC BY-NC-SA 4.0 (https://creativecommons.org/licenses/by-nc-sa/4.0/)
%
%%%%%%%%%%%%%%%%%%%%%%%%%%%%%%%%%%%%%%%%%

%----------------------------------------------------------------------------------------
%	PACKAGES AND OTHER DOCUMENT CONFIGURATIONS
%----------------------------------------------------------------------------------------

\documentclass[
	%a4paper, % Uncomment for A4 paper size (default is US letter)
	10pt, % Default font size, can use 10pt, 11pt or 12pt
]{resume} % Use the resume class
\usepackage[symbol, flushmargin]{footmisc}
\usepackage{utopia} % Use the EB Garamond font
\usepackage{fancyhdr}
\usepackage{lastpage}
\usepackage[danish]{babel}
\usepackage{hyperref}
\hypersetup{
    colorlinks=true,
    linkcolor=black,
    filecolor=magenta,      
    urlcolor=blue
    }
\pagestyle{fancy}
\cfoot{Side \thepage\ af \pageref{LastPage}}
\renewcommand{\headrulewidth}{0pt}

%------------------------------------------------

\name{Kristian Grünberg Olesen Larsen} % Your name to appear at the top

% You can use the \address command up to 3 times for 3 different addresses or pieces of contact information
% Any new lines (\\) you use in the \address commands will be converted to symbols, so each address will appear as a single line.

\address{kristianuruplarsen@gmail.com \\ (+45) 22 93 02 04 } % Contact information
\address{Valborg Allé 19, 5.tv. \\ København, Danmark} % Main address
\address{\href{https://www.linkedin.com/in/kristian-gr\%C3\%BCnberg-olesen-larsen-826a05a8/}{LinkedIn Profil} \\ \href{https://kristianolesenlarsen.github.io/}{Hjemmeside}}
%----------------------------------------------------------------------------------------

\begin{document}

%----------------------------------------------------------------------------------------
%	EDUCATION SECTION
%----------------------------------------------------------------------------------------

\begin{rSection}{Uddannelse}
    \textbf{Københavns Universitet}, \textit{PhD i Økonomi} \hfill \textit{2019-2023} \\
    Vejledere: Søren Leth-Petersen \& Claus Thustrup Kreiner \\
    Afhandling: {Gender Inequality, Labor Supply and New Structural Methods}

    \textbf{Københavns Universitet}, \textit{Cand.polit.} \hfill \textit{2018-2021\footnote{Jeg var indskrevet på PhD-skolen på en ``4+4 kontrakt'', hvorfor jeg var både kandidat- og PhD-studerende indtil halvejs igennem PhD-forløbet.}} \\
    Speciale: {Child Penalties in Labor Returns? Evidence from Linked Survey and Administrative Data}

    \textbf{Københavns Universitet}, \textit{Ba.polit.} \hfill \textit{2014-2017} \\
    Bacheloropgave: {Does Displacement Effects Persist to the Next Generation? Evidence from Danish Childrens Performance in Primary School Exams}
\end{rSection}

%----------------------------------------------------------------------------------------
%	WORK EXPERIENCE SECTION
%----------------------------------------------------------------------------------------

\begin{rSection}{Erfaring}
    \begin{rSubsection}{Skatteministeriets Departement}{2024-nu}{Økonom}{}
        \item[] 
    \end{rSubsection}

    \begin{rSubsection}{CEBI, Økonomisk Institut, Københavns Universitet}{2023-2024}{Postdoctoral Researcher}{}
        \item[] Empirisk og computationel forskning, publikation af forskningsprojekter i videnskabelige tidsskrifter.
    \end{rSubsection}

    \begin{rSubsection}{CEBI, Økonomisk Institut, Københavns Universitet}{2019-2023}{PhD Stipendiat}{}
        \item[] Empirisk og computationel forskning, universitetsundervisning, forskningsformidling. Jeg arbejdede med kollegaer fra København og Princeton University på et forskningsprojekt om optimal beskatning, skrev et solo-papir med emne inden for arbejdsmarkedsøkonomi og anvendte moderne metoder fra reinforcement learning til forskning i computationel økonomi. Jeg havde dagligt ansvar for 4-8 forskningsassistenter (studerende).
    \end{rSubsection}

    \begin{rSubsection}{CEBI, Økonomisk Institut, Københavns Universitet}{2017-2019}{Forskningsassistent}{}
        \item[] Empirisk forskningsassistance bestående I analyser på baggrund af Danmarks Statistiks befolkningsregistre med brug af Stata, SAS og R. Jeg arbejdede primært på sundhedsøkonomiske projekter der har resulteret i publikationer i PNAS og ReStat.
    \end{rSubsection}

    %    \textbf{Københavns Universitet, Økonomisk Institut}\hfill \textit{2023 - Nu} \\
    %    Postdoctoral Researcher
    %
    %    \textbf{Københavns Universitet, Økonomisk Institut}\hfill \textit{2019 - 2023} \\
    %    PhD Stipendiat
    %
    %    \textbf{Københavns Universitet, Økonomisk Institut}\hfill \textit{2017 - 2019} \\
    %    Forskningsassistent ved \textit{Center for Economic Behavior and Inequality (CEBI)}.

\end{rSection}

\begin{rSection}{Undervisningserfaring}

    \textbf{Seminar: Applications of Machine Learning in Economics (MSc)}\hfill \textit{Forår 2021} \\
    Kursuskoordinator og vejleder, Ka. niveau, Økonomisk Institut.

    \textbf{Social Data Science: Econometrics and Machine Learning}\hfill \textit{Forår 2020} \\
    Kursusudvikling \& undervisningsassistent, Ka./PhD niveau, SODAS.

    \textbf{Social Data Science}\hfill \textit{August 2019} \\
    Kursusdrift og undervisningsassistent, Ba./Ka. niveau, SODAS.

    \textbf{Topics in Social Data Science}\hfill \textit{Efterår 2018} \\
    Undervisningsassistent, Ba./Ka./PhD. niveau, SODAS.

    \textbf{Social Data Science}\hfill \textit{August 2018} \\
    Undervisningsassistent, Ba./Ka. niveau, SODAS.
\end{rSection}


\newpage

%----------------------------------------------------------------------------------------
%	Research
%----------------------------------------------------------------------------------------

\begin{rSection}{Forskning}

    \textbf{Micro vs Macro Labor Supply Elasticities: The Role of Dynamic Returns to Effort}\hfill \textit{[\href{https://www.nber.org/papers/w31549}{Working Paper}]} \\
    med Henrik Kleven, Claus Thustrup Kreiner \& Jakob Egholt Søgaard \\
    \textit{R\&R American Economic Review}

    \textbf{Couples and Gender Inequality}\hfill \textit{[\href{https://ssrn.com/abstract=4697847}{Working Paper}]} \\
    Solo papir

    \textbf{Using Reinforcement Learning for Solving Dynamic Discrete-time Problems in Economics} \\
    med Joachim Kahr Rasmussen, \textit{working paper kan fremsendes på forespørgsel}
\end{rSection}


%----------------------------------------------------------------------------------------
%	Services
%----------------------------------------------------------------------------------------

\begin{rSection}{Frivillighed}

    \textbf{Dansk Bjerg og Klatreklub (DBKK)}, \textit{Bestyrelsesmedlem} \hfill \textit{2024} \\
    \textbf{Studenterrådet ved Københavns Universitet}, \textit{Bestyrelsesmedlem} \hfill \textit{2016-2018} \\
    \textbf{Københavns Universitets Uddannelsesstrategiske Råd (KUUR)}, \textit{Studenterrepræsentant} \hfill \textit{2017} \\
    \textbf{Akademisk Råd SAMF}, \textit{Studenterrepræsentant} \hfill \textit{2017} \\

\end{rSection}

%----------------------------------------------------------------------------------------
%	TECHNICAL STRENGTHS SECTION
%----------------------------------------------------------------------------------------

\begin{rSection}{Diverse}

    \begin{tabular}{@{} >{\bfseries}l @{\hspace{6ex}} l @{}}
        Programmeringssprog & Python, R, STATA, SAS (kendskab til SQL, Julia, Matlab) \\
        Talte sprog         & Flydende dansk og engelsk                               \\
    \end{tabular}

\end{rSection}

%----------------------------------------------------------------------------------------
%	EXAMPLE SECTION
%----------------------------------------------------------------------------------------

%\begin{rSection}{Section Name}

%Section content\ldots

%\end{rSection}

%----------------------------------------------------------------------------------------

\end{document}
